\chapterimage{onions.jpg}


\chapter{Una pila de capas}\label{sec:stack_layers}

La comunicación entre ordenadores moderna se consigue mediante una
\concept{stack} de \conceptRef{layer}{capas} de red.
Cada pila se centra en algunos problemas y es responsable de
proporcionar ciertas funcionalidades, \eg,
\conceptRef{address}{direccionamiento} y \concept{sequencing} (secuenciación) correcta (introducidos en los capítulos anteriores).

Cada \concept{layer} define un \textit{par} de funciones \inlineCode{send}/\inlineCode{receive},
que es responsable de algunas de estas funcionalidades.
% 
Cada capa utiliza el \inlineCode{send} y el \inlineCode{receive} de la capa inferior,
\ie, las funcionalidades de una capa se apoyan en las funcionalidades de todas las capas que tiene por debajo.

Para coordinar las capas y hacer posible la comunicación \concept{virtual} entre ellas,
cada capa define un \concept{packet format} (formato de paquete) distinto, que \conceptRef{encapsulation}{encapsula}
un paquete de la capa que tiene por encima.


\section{Envío de datos}

Para implementar estas
capacidades, cada capa transmite no solo los \inlineCode{data} proporcionados por la capa superior,
sino también una cabecera con toda la metainformación necesaria para esa capacidad. Así, cada capa necesita definir
el \concept{packet format}, añadiendo cierto \concept{overhead} (sobrecarga) (típicamente una \conceptRef{header}{cabecera}) al \concept{payload} (carga útil), \eg:

\begin{center}
    \begin{bytefield}{36}
     \bitbox{12}{Header} & \bitbox{24}{Payload}
    \end{bytefield}
\end{center}



Cuando un programa de usuario quiere enviar algunos datos, utiliza el método \inlineCode{send} de la capa \textit{superior} de la pila,
lo que acaba provocando que la capa física transmita ciertos datos, \eg, a través de un cable o una antena.
La siguiente figura representa la estructura de una \concept{stack} de 4 capas:

\begin{center}
\includegraphics[width=\linewidth]{encapsulation_send.pdf}
\end{center}


Los paquetes producidos por cada capa se denominan genéricamente \concept{PDU}s (Protocol Data Unit).
Las PDU de cada capa tienen nombres específicos como \concept{frame} (trama), \conceptRef{datagram}{datagrama} o \conceptRef{segment}{segmento}.
Estos nombres deben usarse con precisión en contextos técnicos.

Cada capa incluye la \concept{PDU} de la capa superior dentro de la definición del \concept{packet format} de su propia PDU.
Este proceso se denomina \conceptRef{encapsulation}{encapsulación}.

\begin{exercise}

El programa de usuario quiere enviar $420$~bytes de datos utilizando la \concept{stack} de la figura anterior.
Supón un tamaño de cabecera para las capas 2, 3 y 4 de $14$~bytes, $20$~bytes y $20$~bytes, respectivamente.
\begin{itemize}
    \item ¿Cuál es el tamaño del \concept{payload} y el tamaño del \concept{overhead} de la \concept{PDU} de cada \conceptRef{layer}{capa}?
    \item ¿Cuál es el tamaño del \concept{payload} y el tamaño del \concept{overhead} de toda la \concept{stack}?
\end{itemize}
\end{exercise}

\section{Recepción de datos}

Cuando un \conceptRef{packet}{paquete} llega a la \concept{stack} de red del destinatario, se procesa en orden inverso,
\ie, las capas se recorren de abajo arriba. Cada capa inspecciona \textit{solo} la \conceptRef{header}{cabecera} correspondiente a esa capa.
Si todo es correcto, la capa elimina la cabecera y pasa el resto del paquete a la capa superior.
A su vez, la capa de arriba pasa los datos restantes (el \concept{payload}) al destinatario.
Continuando con el ejemplo de la pila de 4 capas:

\begin{center}
\includegraphics[width=\linewidth]{encapsulation_receive.pdf}
\end{center}


En una \conceptRef{stack}{pila de red}, solo la capa inferior interactúa directamente con el exterior,
\eg, solo la tarjeta de red física de un ordenador está conectada a un cable que lleva datos a otros ordenadores.
% 
Al mismo tiempo, los datos enviados por la capa inferior incluyen las \conceptRef{PDU}{PDU} de todas las capas
mediante la \concept{encapsulation}.
% 
Gracias a esto, existe una comunicación \concept{virtual} entre cada capa de un extremo y su homóloga en el otro
extremo de la comunicación. La siguiente figura representa la comunicación real (línea continua) y la virtual (línea discontinua)
en nuestro ejemplo:

\begin{center}
    \includegraphics[width=0.5\linewidth]{horizontal_layer_communication.pdf}
\end{center}


\begin{exercise}
Otro ejemplo de \concept{stack} de red tiene 3 capas: L1, L2 y L3 de abajo arriba.

La \concept{PDU} de la capa L2 es la siguiente:

\begin{center}
\begin{bytefield}{32}
\bitheader{0,7,8,15,16,23,24,31}\\
\bitbox{16}{Destination address} & \bitbox{16}{Payload length (bytes)} \\
\wordbox{2}{{Payload (variable length)\\$\vdots$}}
\end{bytefield}
\end{center}

La \concept{PDU} de la capa L3 es la siguiente:

\begin{center}
\begin{bytefield}{32}
\bitheader{0,7,8,15,16,23,24,31}\\
\bitbox{16}{Fragment offset} & \bitbox{16}{Payload length (bytes)} \\
\wordbox{2}{{Payload (variable length)\\$\vdots$}}
\end{bytefield}
\end{center}

La capa L1 recibe un \conceptRef{packet}{paquete} de datos y su función \inlineCode{receive} devuelve los siguientes datos
(supón enteros \concept{unsigned} (sin signo) y orden \concept{big endian} cuando sea necesario:
\begin{center}
\otherBase{01 23 00 0a 01 00 00 06 00 aa 00 bb 43 20}
\end{center}
\begin{itemize}
\item ¿Qué \concept{payload} de datos se envió al programa destinatario en este \conceptRef{packet}{paquete}?\\[-0.25cm]

\item Si este es el último \conceptRef{packet}{paquete} de la comunicación, ¿cuántos datos (de \concept{payload}) se enviaron
  al programa de usuario destinatario?\\[-0.25cm]
  
\item Como máximo, ¿cuántos ordenadores distintos pueden conectarse usando esta \concept{stack}?\\[-0.25cm]

\item Enviamos un fichero como payload de esta pila. ¿Cuál es su tamaño máximo posible?
\end{itemize}
\label{ex:how_it_is_done:three_layer_network}
\end{exercise}

\begin{center}
\showCode{snippets/encapsulation.py}
\end{center}

\begin{exercise}
El código anterior toma unos datos y escribe los bytes que se enviarían a la capa 1
(la \conceptRef{physical layer}{capa física}) de la \conceptRef{network stack}{pila de red} del ejercicio~\ref{ex:how_it_is_done:three_layer_network}.
Implementa el extremo receptor de la pila de la siguiente manera:

\begin{itemize}
\item Implementa una función \inlineCode{layer2_decapsulate(data: bytes) -> (int, bytes)} que tome una PDU de la capa 2
  y devuelva la tupla \inlineCode{(address, payload)}, donde \inlineCode{address} es la dirección de destino y
  \inlineCode{payload} es la PDU de la capa 3 encapsulada.\\[-0.25cm]
  
\item Implementa una función \inlineCode{layer3_decapsulate(data: bytes) -> (int, bytes)} que tome una PDU de la capa 3
  y devuelva la tupla \inlineCode{(offset, payload)}, donde \inlineCode{offset} es el offset (desplazamiento) de la carga útil, y
  \inlineCode{payload} es el fragmento de datos enviado por el otro extremo de la comunicación.\\[-0.25cm]
  
\item Implementa una función \inlineCode{stack_receive(data: bytes)} que decodifique una PDU de la capa 1 (usando las dos funciones anteriores)
  y muestre una línea de mensaje similar a
  \begin{center}
  \inlineCode{"Received a chunk of data. Length: 12 bytes, Address: 0xabcd, Offset: 0."}
  \end{center}\ \\[-0.75cm]
  
\item Completa la función \inlineCode{stack_send} de forma que los datos de entrada se dividan en fragmentos de como mucho 100 bytes,
      se produzca una PDU de la capa 1 válida para cada uno de esos fragmentos,
      y esas PDU de la capa 1 se pasen a \inlineCode{stack_receive}.\\[-0.25cm]
      
\item ¿Funcionaría tu implementación si las PDU de la capa 1 se ordenaran aleatoriamente antes de pasarlas a la pila receptora?
\end{itemize}
% 
\label{ex:how_it_is_done:implementation_three_layer_network}
\end{exercise}

\begin{remark}
La biblioteca \inlineCode{struct} de Python es útil, aunque no obligatoria, para dar formato a los bytes de un \conceptRef{packet}{paquete}.
\end{remark}

